\documentclass[11pt]{article}

\usepackage[margin=1in]{geometry}
\usepackage{amsmath, amssymb}
\usepackage{algorithm}
\usepackage{algpseudocode}
\usepackage{hyperref}

\newcommand{\code}[1]{\texttt{#1}}

\title{CEV Planner Notes}
\author{Patrick Andre}

\date{\today}

\begin{document}
\maketitle

\section{What we're doing}

We have a car that needs to get from the start of the road to the end of it.
\code{rrt.py} and \code{sst.py} load that world with \code{cev\_sim.configure} for the
\code{first-demo} environment and the \code{merceves} vehicle. Positions are meters on the
ground plane, $(x, z)$.

So far we have a basic RRT working, and the car can drive along the path it finds. We've also started
on SST, which plans with the car's actual motion instead of treating it as a point.

The code is split into three files:
\begin{itemize}
  \item \code{world.py} is the old circle map. The runners still use its car helpers; the map comes
    from \code{cev\_sim}.
  \item \code{rrt.py} is the RRT planner plus the code that makes the car follow the path.
    Run it with \code{python rrt.py}.
  \item \code{sst.py} is the SST work so far. It doesn't plan yet; running it just lets you drive the car
    around with WASD.
\end{itemize}
You need the \code{cev\_sim} package from the simulator checkout, plus \code{numpy} and
\code{matplotlib}. The map window is the library's planner view.

\section{The car}

The car's position is the middle of its rear axle. It has a position $(x, y)$, heading $\theta$, speed
$v$, and turning rate $\omega$. We control it with an acceleration $u_a$ and a steering rate $u_\omega$
(how fast the turning rate changes):
\[
  \dot x = v\cos\theta, \quad \dot y = v\sin\theta, \quad \dot\theta = \omega, \quad
  \dot v = u_a, \quad \dot\omega = u_\omega.
\]
\code{rrt.py} still computes that heading rate. \code{cev\_sim} takes a speed fraction and a steering
fraction, both in $[-1, 1]$, and the heading rate is turned into the steering fraction before \code{step}.

The important thing here is that the car can't change how sharply it's turning instantly, since we only
control how fast the turning rate changes. That ends up mattering a lot when we try to follow RRT's path.

\section{Collision checking}

A point hits an obstacle if it's within the obstacle's radius plus a small margin (5 px, half the car's
width) of its centre. Going off the edge of the map also counts as a collision.

To check a whole line segment, we walk along it and check a point every 2 px or so. Only checking the
end of a segment isn't enough, because a 20 px step can jump right over one of the small green obstacles.

Checking every point against all 377 obstacles is slow, so we sort the obstacles by their $x$ coordinate
once at the start. For a segment, an obstacle can only matter if its centre's $x$ is within 79 px of the
segment (79 = the biggest radius, 74, plus the 5 px margin). We use binary search (\code{bisect}) to grab
just the obstacles in that range. That cuts it down to around 32 obstacles per check instead of 377, and
in our test it was about 14 times faster with exactly the same results.

\section{RRT}

RRT builds a tree starting from the start point. Each round we pick a random point on the map, find the
closest node already in the tree, and take a small step from that node towards the random point. If the
step doesn't hit anything, it becomes a new node. Since the random points are spread all over the map,
the tree ends up spreading out quickly in every direction.

For now we treat the car as a point that can move in any direction, so the tree is made of straight lines.
Each node just remembers its position and which node it came from (its parent). The tree is a regular
Python list of nodes, and the parent links are what actually make it a tree.

\begin{algorithm}[h]
\caption{Our RRT}
\begin{algorithmic}[1]
\State tree $\gets$ [start]
\For{up to 5000 rounds}
  \State $q_{\text{rand}} \gets$ random point (or the target 5\% of the time)
  \State $q_{\text{near}} \gets$ closest node in tree to $q_{\text{rand}}$
  \State $q_{\text{new}} \gets$ step 20 px from $q_{\text{near}}$ towards $q_{\text{rand}}$
  \If{the line from $q_{\text{near}}$ to $q_{\text{new}}$ is collision free}
    \State add $q_{\text{new}}$ to the tree with parent $q_{\text{near}}$
    \If{$q_{\text{new}}$ is inside the target}
      \State \Return the path from start to $q_{\text{new}}$
    \EndIf
  \EndIf
\EndFor
\State \Return no path found
\end{algorithmic}
\end{algorithm}

A few details:

\paragraph{Goal bias.} 5\% of the time, instead of a random point we use the target itself. Without this
the tree eventually gets there, but random points almost never land in the target so it takes a while.
If you set it too high the tree gets greedy and keeps running into obstacles between it and the goal.

\paragraph{Taking a step.} To step from $q_{\text{near}}$ towards $q_{\text{rand}}$ we get the direction
with $\phi = \operatorname{atan2}(\Delta y, \Delta x)$ and move 20 px that way:
\[
  q_{\text{new}} = q_{\text{near}} + 20\,(\cos\phi, \sin\phi).
\]
If $q_{\text{rand}}$ is already closer than 20 px we just go straight to it. Keeping the steps small means
the tree grows gradually instead of drawing huge lines across the whole map.

\paragraph{Getting the path.} Once a node lands in the target, we follow the parent links back to the
start and flip the list around so it goes start to goal.

\paragraph{How it does.} We ran it 30 times with different random seeds and it found the goal in 29 of
them. A typical run builds around 500 nodes, ends up with a path of about 128 points, and takes about a
tenth of a second. You get a different tree every time since it's random.

\section{Making the car follow the path}

RRT just gives us a list of points, so we still need to actually drive the car along them. Every frame
the controller:
\begin{enumerate}
  \item Picks a point on the path about 25 px ahead of the car to aim at.
  \item Works out how far off the car's heading is from that point ($e_\theta$), and turns towards it:
    the bigger the error, the harder it turns ($\omega_{\text{want}} = 4\,e_\theta$, and the steering rate
    pushes the actual turning rate towards that).
  \item Drives at about 80 px/s, slows down when it has to turn sharply, and stops once it's inside the
    target.
\end{enumerate}
If you hold W, A, S or D you take over, and when you let go it goes back to following the path.

\paragraph{The problem.} The car makes it to the goal every time (about 40 seconds), but it cuts corners
and clips obstacles by up to about 11 px on sharp turns. RRT's path is a zig-zag of straight lines that
hug the obstacles, and the car can't turn on the spot. We tried tuning the controller to follow the path
more tightly, but the car just started wobbling and it got worse. We also tried giving the path more space
from the obstacles, but at 8 px of margin or more RRT couldn't find any path at all because the gaps are
too narrow. So this is really a limit of planning for a point and then driving a car along it, which is a
big part of why we want SST.

\section{SST (still in progress)}

SST stands for Stable Sparse RRT (Li, Littlefield and Bekris, 2016). The big difference from what we
have is that instead of straight lines, every branch of the tree comes from actually simulating the car
with some random controls. That means any path it finds is one the car can really drive.

It also changes RRT in two ways:
\begin{itemize}
  \item When picking which node to grow from, it looks at all the nodes near the random sample and picks
    the one with the lowest cost so far, instead of just the closest one.
  \item It keeps the tree small by only keeping the best node in each small region and pruning the rest.
\end{itemize}

\subsection{State and controls}
For SST we're using a simpler car model with state $(x_1, x_2, a, s)$, meaning position, heading and
speed. The controls are a turning rate $u_a$ and an acceleration $u_s$:
\[
  \dot x_1 = s\cos a, \quad \dot x_2 = s\sin a, \quad \dot a = u_a, \quad \dot s = u_s.
\]
This is slightly different from the \code{Car} class, where the turning rate is part of the state rather
than something we control directly. We'll have to sort that out before we can drive an SST path with
\code{Car}.

To measure how close two states are, we push each one forward along its heading for one second
($\bar x_1 = x_1 + s\cos a$, same for $\bar x_2$ with $\sin$) and take the distance between those points.
That way two states count as close if they're near each other \emph{and} heading the same way.

\subsection{Monte Carlo propagation}
This is how SST grows a branch: pick a random amount of time $t$ between 0 and 0.5 s, pick random
controls, and work out where the car ends up:
\[
  x_{\text{new}} = x_{\text{prop}} + \int_0^t f(x(\tau), u)\, d\tau.
\]
Since the controls stay constant over the step, we can do the integral by hand (integration by parts) and
get a closed form. With $a_0$ and $s_0$ as the starting heading and speed, and $a_1 = a_0 + u_a t$:
\begin{align*}
  x_1 &= x_1^0 + \frac{s_0}{u_a}(\sin a_1 - \sin a_0) + \frac{u_s}{u_a^2}(\cos a_1 - \cos a_0 + u_a t \sin a_1), \\
  x_2 &= x_2^0 - \frac{s_0}{u_a}(\cos a_1 - \cos a_0) + \frac{u_s}{u_a^2}(\sin a_1 - \sin a_0 - u_a t \cos a_1),
\end{align*}
and the new heading and speed are just $a_1$ and $s_0 + u_s t$. We checked these against a numerical
integration and they match. One catch is that this divides by $u_a$, so it breaks if the car is going
perfectly straight ($u_a = 0$). We'll deal with that if it comes up.

\subsection{Picking a node}
\code{near} finds all the nodes within 7 of a random state, plus the closest one overall.
\code{bestFirstSelectionSST} then picks the cheapest node out of the nearby ones, or the closest node if
nothing is nearby.

\subsection{What's left}
\begin{itemize}
  \item The main SST loop (pick a node, propagate, check for collisions, add it).
  \item The pruning part, and deciding what ``cost'' actually is (probably time).
  \item Collision checking along curved paths, since \code{collision\_free} only does straight lines.
  \item \code{getRandomState} and \code{getRandomAction} still call \code{math.pi()}, which should be
    \code{math.pi}.
  \item Better ranges for the random controls (turning both ways, braking as well as speeding up).
  \item Switching \code{SSTTree} to plain lists instead of NumPy arrays, since we append to it a lot.
\end{itemize}

\end{document}
